% ===========================================================================
%  Entropy -- a Loom fill-in notebook.
%  Original exposition of standard (public-domain) information theory.
%  Compile:  latexmk -xelatex main      (or: xelatex main, twice)
% ===========================================================================
\documentclass{loom}

\newcommand{\E}{\mathbb{E}}
\newcommand{\R}{\mathbb{R}}
\newcommand{\X}{\mathcal{X}}
\newcommand{\Ent}{\mathrm{H}}
\newcommand{\MI}{\mathrm{I}}
\newcommand{\KL}{\mathrm{D}}
\newcommand{\abs}[1]{\lvert #1\rvert}

\runningthread{entropy}

\begin{document}

\loomcover
  {Entropy}
  {The price of uncertainty · fill-in notes}
  {woven on \textsc{Loom}}
  {\today}

% ===========================================================================
\section*{The one idea, and how to read these notes}
\addcontentsline{toc}{section}{The one idea · how to read}
\warmth{0}

\begin{strand}
\keyword{Entropy} $\Ent(X)$ is the average \emph{surprise} of a random source, measured in bits.
The single idea of these notes: entropy is a \keyword{price}. It is the exact number of bits to
describe, store, or transmit $X$, and every theorem in the subject is one of two sentences:
\keyword{you cannot pay less} (a converse) or \keyword{you needn't pay more} (an achievability).
\end{strand}

\block{The master table: the price, in bits (fill it in as you go)}
\noindent\begin{tabularx}{\linewidth}{@{}l L@{}}
\toprule
{\color{indigo}quantity} & \multicolumn{1}{l}{\color{madder}what you are paying for}\\
\midrule
$\Ent(X)$               & bits to describe $X$ $=$ the best \fillin[3cm] rate \;(§\ref{sec:source})\\
$\Ent(X\mid Y)$         & bits still needed for $X$ once you know $Y$\\
$\MI(X;Y)$              & bits $Y$ tells you about $X$ $=\Ent(X)-\fillin[2.5cm]$ \;(§\ref{sec:toolkit})\\
$\KL(p\,\Vert\,q)$      & extra bits for coding $p$ with the wrong model $q$ \;($\ge\fillin[0.8cm]$)\\
channel capacity $C$    & $\max_p \MI(X;Y)$: bits you can send reliably \;(§\ref{sec:faces})\\
\bottomrule
\end{tabularx}

\block{How to read (passive \emph{and} active)}
\begin{itemize}[leftmargin=1.3em,itemsep=1pt]
  \item \keyword{Read} the strands and the stated theorems.
  \item \keyword{Fill} the blanks \fillin[1.2cm], the \TODO{missing step} proof gaps, and the
        \emph{Your turn} boxes. Any information-theory text is the answer key.
  \item Bump each \texttt{\textbackslash warmth\{0\}} to $0$--$5$ once you can rebuild it cold.
  \item Logs are base $2$ (bits) unless noted; base $e$ gives nats.
\end{itemize}

\bigskip
{\small\tableofcontents}

\clearpage
% ===========================================================================
\section{Surprise, and its average}\label{sec:entropy}
\warmth{0}\quad\whisper{Rare events are surprising; entropy is how surprised you are on average.}

\begin{strand}
Quantify \keyword{surprise}: an event of probability $p$ carries $-\log p$ bits of it. Certain
events ($p=1$) surprise you not at all; rare events ($p\to0$) surprise you a lot; and surprises
of \emph{independent} events add. Entropy is just the \emph{expected} surprise of a source: how
many yes/no questions, on average, to pin down its value.
\end{strand}

\block{The definition}

\begin{definition}[entropy]\label{def:entropy}
For a discrete random variable $X$ on alphabet $\X$ with law $p$,
\[
  \Ent(X)\;=\;\E\big[-\log p(X)\big]\;=\;-\sum_{x\in\X} p(x)\,\fillin[2cm].
\]
It depends only on the \emph{probabilities}, not on the values $x$. In bits ($\log_2$) it is the
average number of fair-coin questions needed to identify $X$.
\end{definition}

\recall{Three sanity checks: a sure thing has $\Ent=0$; a fair coin has $\Ent=1$ bit; and $\Ent$ is maximized by the \emph{uniform} law (most uncertain), where it equals $\log\abs{\X}$ (§\ref{sec:source}).}

\begin{yourturn}
Compute (in bits):
\[
  \Ent(\text{fair coin})=\fillin[1cm],\qquad
  \Ent\big(\tfrac14,\tfrac34\big)=\fillin[1.6cm],\qquad
  \Ent(\text{uniform on } 8)=\fillin[1cm].
\]
\recall{$1$ bit; $\ -\tfrac14\log\tfrac14-\tfrac34\log\tfrac34\approx 0.811$ bits; $\ \log_2 8=3$ bits. A biased coin is \emph{less} uncertain than a fair one, so it is cheaper to describe.}
\end{yourturn}

\begin{strand}
Why $-\log p$ and not some other measure of surprise? Because it is the \emph{only} one (up to the
choice of base) that is continuous, that grows as events get rarer, and that makes the surprise of
independent events \keyword{add}: $-\log(pq)=-\log p-\log q$. Additivity over independence is the
hinge the whole subject swings on.
\end{strand}

\loose{Continuous $X$: replace the sum by an integral to get \emph{differential} entropy $h(X)=-\int f\log f$. It can be negative and is not coordinate-free, so handle with care; the clean, operational quantity there is relative entropy (§\ref{sec:toolkit}).}

% ===========================================================================
\section{The toolkit}\label{sec:toolkit}
\warmth{0}\quad\whisper{Joint, conditional, shared, and wasted bits — all bookkeeping on one quantity.}

\begin{strand}
From the single notion of entropy, four relatives organize everything: \keyword{joint} entropy
(two sources together), \keyword{conditional} entropy (what's left after seeing one),
\keyword{mutual information} (what they share), and \keyword{relative entropy} (the cost of a wrong
model). They are tied by one picture and a couple of identities.
\end{strand}

\block{The four quantities}

\begin{definition}[the relatives]\label{def:relatives}
$\displaystyle \Ent(X,Y)=\E[-\log p(X,Y)]$, $\quad\Ent(X\mid Y)=\E[-\log p(X\mid Y)]$,
\[
  \MI(X;Y)=\Ent(X)-\Ent(X\mid Y),\qquad
  \KL(p\,\Vert\,q)=\sum_x p(x)\log\frac{p(x)}{q(x)} .
\]
\end{definition}

\block{The identities (the information diagram)}
\noindent\begin{tabularx}{\linewidth}{@{}l L@{}}
\toprule
{\color{indigo}identity} & \multicolumn{1}{l}{\color{indigo}reading}\\
\midrule
$\Ent(X,Y)=\Ent(X)+\Ent(Y\mid X)$            & \keyword{chain rule}: total $=$ first $+$ rest\\
$\MI(X;Y)=\Ent(X)+\Ent(Y)-\Ent(X,Y)$         & shared $=$ overlap of the two circles\\
$\MI(X;Y)=\KL\!\big(p_{XY}\,\Vert\,p_X p_Y\big)$ & information $=$ distance from \keyword{independence}\\
$\MI(X;Y)\ge 0$, with $=0$ iff independent    & you never lose by looking at $Y$\\
\bottomrule
\end{tabularx}

\begin{strand}
Picture two overlapping circles, $\Ent(X)$ and $\Ent(Y)$. The overlap is $\MI(X;Y)$; the
left-only crescent is $\Ent(X\mid Y)$; their union is $\Ent(X,Y)$. Every identity above is just a
statement about areas in this Venn diagram, which is why information theorists draw it constantly.
\end{strand}

\begin{yourturn}
Fill the relations from the diagram:
\[ \MI(X;Y)=\Ent(Y)-\fillin[2.5cm]. \]
\[ \Ent(X\mid Y)\ \fillin[0.8cm]\ \Ent(X)\qquad(\text{conditioning can only \fillin[2.5cm]}). \]
\recall{$\MI(X;Y)=\Ent(Y)-\Ent(Y\mid X)$ (symmetric!); and $\Ent(X\mid Y)\le\Ent(X)$: conditioning can only \emph{reduce} entropy on average. Knowing more never hurts.}
\end{yourturn}

\begin{strand}
A caution: $\KL$ is \emph{not} a metric (not symmetric, no triangle inequality). Yet it is the
right notion of ``distance to the truth'': the per-symbol penalty in bits for compressing data
from $p$ as if it came from $q$. That operational meaning is §\ref{sec:source}.
\end{strand}

% ===========================================================================
\section{Two pillars: $\KL\ge0$ and the compression limit}\label{sec:source}
\warmth{0}\quad\whisper{One inequality implies all the bounds; one theorem turns entropy into bits on a disk.}

\begin{strand}
The whole edifice rests on two results. The first is an \emph{inequality} from which every other
bound drops out: relative entropy is never negative. The second is the \emph{operational} payoff
that makes the bits real: Shannon's source coding theorem, which says the average code length can
get down to the entropy, and no further.
\end{strand}

\block{Pillar 1: nothing beats the truth}

\begin{theorem}[Gibbs' inequality]\label{thm:gibbs}
$\KL(p\,\Vert\,q)\ge 0$ for all distributions $p,q$, with equality iff $p=q$.
\end{theorem}
\begin{proof}
Apply Jensen to the concave $\log$:
$-\KL(p\Vert q)=\sum_x p(x)\log\frac{q(x)}{p(x)}\le\log\sum_x p(x)\frac{q(x)}{p(x)}=\log 1=0$.
\TODO{state the equality condition from strict concavity, and derive the two corollaries below}
\end{proof}

\recall{Two corollaries, for free, by choosing $q$: take $q$ uniform to get $\Ent(X)\le\log\abs{\X}$ (uniform is maximally uncertain); take $q=p_Xp_Y$ to get $\MI(X;Y)\ge0$, i.e.\ $\Ent(X\mid Y)\le\Ent(X)$.}

\block{Pillar 2: entropy is the floor on code length}

\begin{theorem}[Kraft + source coding]\label{thm:source}
Any binary \keyword{prefix} code with lengths $\ell(x)$ obeys $\sum_x 2^{-\ell(x)}\le 1$ (Kraft).
Consequently the expected length $L=\sum_x p(x)\ell(x)$ satisfies
\[
  L\;\ge\;\Ent(X),
\]
and a code achieving $L<\Ent(X)+1$ always exists (take $\ell(x)=\lceil-\log p(x)\rceil$).
\end{theorem}
\begin{proof}
For the lower bound, $L-\Ent(X)=\sum_x p(x)\log\frac{p(x)}{2^{-\ell(x)}}=\KL(p\Vert r)+\log\frac1{\sum 2^{-\ell}}\ge0$,
where $r(x)=2^{-\ell(x)}/\sum 2^{-\ell}$. \TODO{check both terms are $\ge0$ using Kraft and Gibbs}
\end{proof}

\begin{yourturn}
Code the distribution $p=\big(\tfrac12,\tfrac14,\tfrac18,\tfrac18\big)$.
\[
  \Ent(X)=\fillin[1.6cm]\text{ bits},\qquad L=\fillin[1.6cm]\text{ bits}.
\]
\[
  \text{optimal code lengths}\ (\ell_1,\ell_2,\ell_3,\ell_4)=\fillin[2.6cm].
\]
\recall{$\Ent=1.75$ bits; optimal lengths $(1,2,3,3)$; $L=\Ent$ exactly because the probabilities are \keyword{dyadic} (powers of two). Otherwise you pay up to one extra bit.}
\end{yourturn}

\begin{strand}
Squeeze out the last bit by coding \emph{blocks}: over $n$ symbols the per-symbol overhead is
$1/n\to0$, so the rate $\to\Ent(X)$. The clean statement is the AEP / typical set: about
$2^{n\Ent}$ sequences carry essentially all the probability, so $n\Ent$ bits suffice (§\ref{sec:faces}).
\end{strand}

% ===========================================================================
\section{Faces of entropy}\label{sec:faces}
\warmth{0}\quad\whisper{Training losses, channel limits, the most honest prior: all are entropy, wearing hats.}

\begin{strand}
The same price tag shows up across machine learning and communication. Each face below is entropy
or one of its relatives, doing a job.
\end{strand}

\block{The keyring}
\noindent\begin{tabularx}{\linewidth}{@{}l L@{}}
\toprule
{\color{indigo}face} & \multicolumn{1}{l}{\color{indigo}the idea}\\
\midrule
\keyword{cross-entropy} loss   & training minimizes $\E[-\log q_\theta]=\Ent(p)+\KL(p\,\Vert\,q_\theta)$, i.e.\ it closes the KL gap to the data\\
channel \keyword{capacity}     & $C=\max_p\MI(X;Y)$; you can communicate reliably at any rate below $C$, and not above\\
\keyword{maximum entropy}      & the least-committed law with given constraints: fixed variance $\Rightarrow$ Gaussian, fixed mean on $\R_{\ge0}\Rightarrow$ exponential\\
\keyword{AEP} / typical set    & $\approx 2^{n\Ent}$ sequences carry almost all probability, so $n\Ent$ bits suffice\\
\keyword{Fano}                 & if you must guess $X$ from $Y$, your error is bounded below by $\Ent(X\mid Y)$\\
\bottomrule
\end{tabularx}

\begin{yourturn}
The cross-entropy you minimize when training a classifier is $\E_{p}[-\log q_\theta(X)]$.
\begin{itemize}[leftmargin=1.3em,itemsep=1pt]
  \item It splits as $\Ent(p)+\fillin[3cm]$, and only the second term depends on $\theta$.
  \item So minimizing cross-entropy is exactly minimizing \fillin[3cm] from the model to the data,
        and the floor it can reach is $\fillin[1.6cm]$.
\end{itemize}
\recall{Splits as $\Ent(p)+\KL(p\Vert q_\theta)$; training minimizes $\KL(p\Vert q_\theta)$; the floor is $\Ent(p)$, the data's own entropy, reached only when $q_\theta=p$.}
\end{yourturn}

\block{The one idea, recovered}
\begin{strand}
Read the master table again. Each row is a number of bits, and each theorem is one of two
sentences. \keyword{Converses} (Gibbs, $L\ge\Ent$, capacity, Fano) say a price cannot be undercut;
\keyword{achievabilities} (Shannon codes, the AEP, capacity-achieving schemes) say it can be met.
Description, compression, transmission, and learning are all the same act of paying entropy, and
not a bit more than you must.
\end{strand}

\loose{Where to go next: drop probability for a single string and you get \keyword{Kolmogorov complexity} (the length of the shortest program that prints it); allow lossy reconstruction and you get rate--distortion (the price of ``good enough''); both are entropy's siblings.}


\end{document}
